\begin{abstract}
Kompleksitas arsitektur microservice menyebabkan proses identifikasi sumber penyebab gangguan (Root Cause Localization) menjadi tantangan, terutama karena setiap service memiliki karakteristik metrik dan pola interaksi yang saling berkaitan. Pendekatan pembelajaran mesin berbasis data berlabel dapat digunakan untuk membantu proses lokalisasi, tetapi ketersediaan label gangguan sering kali terbatas. Salah satu pendekatan yang berpotensi mengatasi keterbatasan tersebut adalah Self-Supervised Learning (SSL), yang memungkinkan model mempelajari representasi data tanpa memerlukan label pada tahap prapelatihan. Penelitian ini bertujuan mengevaluasi penerapan SSL dan Transfer Learning untuk lokalisasi service penyebab gangguan pada sistem microservice, dengan membandingkan tiga pendekatan, yaitu TS2Vec, TimeMAE, dan RUN. Evaluasi dilakukan menggunakan dataset RCAEval, yang mencakup tiga sistem microservice, yaitu Online Boutique, Sock Shop, dan Train Ticket. Penelitian ini dirancang melalui dua skenario eksperimen, yaitu evaluasi dalam dataset yang sama (in-domain) dan transfer representasi antar-dataset (cross-dataset transfer). Pada tahap prapelatihan, model mempelajari pola data deret waktu tanpa menggunakan label, kemudian representasi yang diperoleh digunakan pada tahap penyesuaian model (fine-tuning) dengan variasi strategi FineLast dan FineAll. Kinerja lokalisasi dievaluasi menggunakan metrik Accuracy at k (AC@k) dan Average at 5 (Avg@5), dengan membandingkan hasil prapelatihan terhadap inisialisasi acak serta pendekatan heuristik. Penelitian ini diharapkan dapat memberikan pemahaman mengenai kemampuan representasi hasil pembelajaran mandiri dalam mendukung lokalisasi akar penyebab gangguan serta efektivitas transfer representasi antar-sistem microservice.


\keywords{Root Cause Localization \and Microservice \and Self-Supervised Learning \and Transfer Learning \and Representasi deret waktu}
\end{abstract}
